\chapter*{Введение} % * не проставляет номер
\addcontentsline{toc}{chapter}{Введение} % вносим в содержание



\subsection*{Актуальность исследования}
Современное сельское хозяйство сталкивается с комплексом взаимосвязанных 
вызовов: истощение природных ресурсов, сокращение пригодных для обработки 
земель и климатическая нестабильность.\cite{Manthena2022, Budhlakoti2022} Особенно остро эти проблемы 
проявляются в регионах рискованного земледелия, где продуктивность 
культур напрямую зависит от устойчивости сортов к засухе, полеганию 
и другим стрессовым факторам. Традиционная селекция, основанная на 
фенотипическом отборе, требует значительных временных затрат и не всегда 
позволяет эффективно выявлять генетические детерминанты хозяйственно 
ценных признаков.\cite{Crossa2017} 

Современные методы геномного прогнозирования открывают возможность 
существенного ускорения селекционного процесса за счёт анализа 
однонуклеотидных полиморфизмов (SNP). Статистические подходы 
(rrBLUP, GBLUP, байесовские методы) обладают ограниченной способностью 
улавливать нелинейные и эпистатические взаимодействия между генами. 
Модели глубокого обучения, в частности архитектуры на основе трансформеров, 
демонстрируют высокую эффективность в задачах выявления сложных 
зависимостей в геномных данных, однако их применение для прогнозирования 
продуктивности сельскохозяйственных культур остаётся недостаточно 
изученным. В связи с этим разработка интерпретируемой модели глубокого 
обучения, адаптированной к специфике генетических данных и способной 
работать с различными культурами, является актуальной задачей.

\subsection*{Объект и предмет исследования}
\textbf{Объект исследования} — процесс прогнозирования фенотипических 
признаков сельскохозяйственных растений по генотипическим данным.

\textbf{Предмет исследования} — архитектура и методы обучения нейросетевой 
модели на основе трансформеров и механизмов внимания для задач геномного 
прогнозирования.

\subsection*{Цель и задачи исследования}
\textbf{Цель работы} — разработать и программно реализовать модель 
глубокого обучения для прогнозирования продуктивности сельскохозяйственных 
растений по SNP-маркерам, обеспечивающую высокую точность предсказания 
и возможность интерпретации выявленных генных взаимодействий.

Для достижения цели поставлены следующие \textbf{задачи}:
\begin{enumerate}
	\item Провести анализ существующих методов геномного прогнозирования, 
	выявить их преимущества и ограничения.
	\item Разработать архитектуру модели, включающую механизм внимания 
	на уровне SNP и реконструкционную регуляризацию.
	\item Реализовать программный комплекс для обучения и тестирования 
	модели на данных различных сельскохозяйственных культур.
	\item Провести вычислительные эксперименты на данных нута, риса и 
	рапса, оценив точность прогнозирования различных фенотипических 
	признаков.
	\item Исследовать стратегии отбора генов и методы аугментации данных 
	для повышения обобщающей способности модели.
	\item Выполнить интерпретацию модели с помощью значений Шепли для 
	выявления значимых межгенных взаимодействий.
\end{enumerate}

\subsection*{Гипотеза исследования}
Предполагается, что введение механизма внимания на уровне отдельных 
SNP-маркеров и реконструкционной регуляризации позволит повысить 
устойчивость обучения и сохранить индивидуальные различия между 
образцами, а разработанная архитектура будет демонстрировать 
сопоставимую эффективность на данных различных культур.

\subsection*{Методологическая база исследования}
Работа базируется на методах машинного обучения и глубокого обучения 
(трансформеры, механизмы внимания, многозадачное обучение), методах 
обработки геномных данных (частотное кодирование SNP, агрегация по 
генам, позиционное кодирование), а также на подходах к интерпретации 
нейросетевых моделей (значения Шепли). Вычислительные эксперименты 
проводились с использованием библиотек PyTorch, NumPy, Pandas, scikit-learn.

\subsection*{Научная новизна}
\begin{enumerate}
	\item Предложена архитектура модели глубокого обучения с механизмом 
	внимания на уровне SNP, позволяющая адаптивно взвешивать вклад 
	отдельных маркеров при формировании генных представлений.
	\item Введена реконструкционная регуляризация скрытых представлений 
	генов, способствующая сохранению информации при прохождении через 
	трансформерный блок.
	\item Разработана многозадачная функция потерь с адаптивным 
	взвешиванием фенотипических признаков на основе обучаемых параметров 
	неопределённости.
	\item Проведено систематическое сравнение стратегий отбора генов 
	(взвешенной, равномерной по геному, равномерной по хромосомам) 
	применительно к трансформерным архитектурам.
\end{enumerate}

\subsection*{Теоретическая и практическая значимость}
\textbf{Теоретическая значимость} заключается в развитии методов 
применения трансформерных архитектур для задач геномного прогнозирования, 
в частности в обосновании эффективности иерархического внимания 
(SNP → ген → геном).

\textbf{Практическая значимость} состоит в создании программной 
реализации модели, которая может быть использована в селекционных 
программах для идентификации генетических маркеров, связанных с 
продуктивностью, устойчивостью к засухе и другими хозяйственно 
ценными признаками. Результаты experiments на трёх культурах 
подтверждают применимость подхода в регионах рискованного земледелия 
для ускоренного выведения адаптированных сортов.


%% Вспомогательные команды - Additional commands
%\newpage % принудительное начало с новой страницы, использовать только в конце раздела
%\clearpage % осуществляется пакетом <<placeins>> в пределах секций
%\newpage\leavevmode\thispagestyle{empty}\newpage % 100 % начало новой строки